\section{Repository Structure and Configuration}

\subsection{The reference implementation and the research branches}

The repository contains a maintained detector, preserved scientific code,
research data and results, supporting explanations, and a separate thesis
workspace. These have different roles. Files existing somewhere in repository
history do not establish that one installed command contains all their
behaviour. The reference detector is \texttt{master} at \texttt{db2a736};
the current writing branch adds the workspace at \texttt{0d66bf8}.
The temporal and multi-signal branch heads remain \texttt{b0aa35d} and
\texttt{4279d84} respectively.

\begin{table}[H]
\centering\small
\begin{tabularx}{\textwidth}{>{\raggedright\arraybackslash}p{0.25\textwidth}X}
\toprule\textbf{Location} & \textbf{Role} \\
\midrule
\path{src/lm_idnet/} & Installed detector, schemas, input processing,
numerical methods, state, CLI, and evaluation. \\
\path{configs/} & Three device research configurations, twelve policy/variant
experiment configurations, and one prepared edge configuration. \\
\path{data/} & Selected processed captures and independent labels, plus the
larger D-Link raw archive. Configured captures are a subset of that archive. \\
\path{artifacts/}, \path{reports/} & Model/threshold pairs, events, adaptation
records, data checks, benchmarks, and simulation measurements. \\
\path{unsw_iot_attack_pcaps/} & Download, symlink preparation, official
annotation metadata, and label preparation. \\
\path{lm-algorithm/} & Preserved authors' numerical capsule and example data;
not the installed application entry point. \\
\path{resouces/} & Supplied brief, research papers, and explanatory material.
The directory spelling is retained for existing references. \\
\path{tests/} & Unit, numerical, subprocess, storage, PCAP, and optional
production-container checks with small fixtures. \\
\path{docs/}, \path{thesis/} & Technical/research records and the independently
buildable thesis skeleton. \\
\bottomrule
\end{tabularx}
\caption{The main repository areas. Detailed file maps appear in the appendices.}
\end{table}

\subsection{Installation and dependencies}

\path{pyproject.toml} is the authoritative Python package definition. It names
Python 3.11 or later, defines the \texttt{lm-idnet} console command, and pins
mpmath, NumPy, pandas, Pydantic, Scapy, SciPy, and tabulate. Pytest is a test
extra rather than a production dependency. The \path{main.py} file is a
compatibility launcher. It does not define a competing pipeline.

\begin{table}[H]
\centering\small
\begin{tabular}{lll}
\toprule
\textbf{Package} & \textbf{Pinned version} & \textbf{Use} \\
\midrule
\texttt{mpmath} & 1.4.1 & Runtime \\
\texttt{numpy} & 2.4.6 & Runtime \\
\texttt{pandas} & 3.0.3 & Runtime \\
\texttt{pydantic} & 2.13.4 & Runtime \\
\texttt{scapy} & 2.7.0 & Runtime \\
\texttt{scipy} & 1.17.1 & Runtime \\
\texttt{tabulate} & 0.10.0 & Runtime \\
\texttt{pytest} & 9.1.1 & Test extra \\
\bottomrule
\end{tabular}
\caption{Declared package versions at the reference snapshot; historical measurements retain their own environment records.}
\end{table}


The native shared library and its coefficient files are packaged with the
application. Docker compiles the library during image creation. Local editable
installation needs the documented compiler step first. A binary built for one
machine must not be assumed usable on another architecture. Older numerical
records name their own Python, NumPy, and SciPy versions; they are not rewritten
to look as though they were measured in today's pinned environment.

\subsection{Research configuration}

An \path{AppConfig} records source and processed roots, dataset and device
identity, the MAC selector, category order, trained window duration, capture
partitions, seed settings, estimator, calibration, adaptation, and outputs.
Pydantic rejects unknown fields. Capture IDs must contain valid dates, remain
unique, and form strictly ordered, disjoint fit, calibration, development, and
final groups. Category count must agree with the estimator, and adaptive buffers
must be large enough for the requested updates.

The current device configurations use four categories, ten-minute windows,
six-digit native LM, raw scores, and the one-percent lower calibration quantile.
The fitting tolerance is $2\times10^{-9}$ with at most 100,000 iterations.
Configuration makes these choices explicit; their presence is not evidence
that every alternative has been experimentally compared.

\subsection{Edge configuration}

An \path{EdgeConfig} retains detector input semantics and learning settings but
does not require raw-data roots or a list of future capture dates. It identifies
a prepared state directory and compatible model/threshold paths. The snapshot
contains the baseline pair, current pair, safety references, and calibration
volume limit, so restoration does not need research captures.

Both configuration types share the same detector fields. The rollback command
accepts either strict format; the run command requires the edge format.
This separation removes a research-only requirement without weakening device,
category, duration, backend, precision, policy, or model--threshold checks.

\subsection{Supporting files and local material}

The root README, Docker ignore rules, Git ignore rules and attributes, and the
commit-instruction location support installation and collaboration. Local
environment and installation-marker files are also present. Their contents are
not thesis evidence and are not reproduced here. Generated logs and
\path{out/} builds are local outputs. Optional raw UNSW staging is ignored and
is not redistributed as though it were part of a small public fixture.
